% Resume: Arumalla Mohit Krishna
% Compile with pdfLaTeX.

\documentclass[letterpaper,10pt]{article}

\usepackage[margin=0.5in]{geometry}
\usepackage[T1]{fontenc}
\usepackage{lmodern}
\usepackage{microtype}
\usepackage{titlesec}
\usepackage{enumitem}
\usepackage[hidelinks]{hyperref}
\usepackage{tabularx}
\input{glyphtounicode}

\pdfgentounicode=1
\pagestyle{empty}
\urlstyle{same}
\raggedright
\setlength{\parindent}{0pt}
\setlength{\tabcolsep}{0pt}

\titleformat{\section}
  {\large\scshape}{}{0em}{}[\titlerule]

\titlespacing*{\section}{0pt}{6pt}{3pt}

\newcommand{\entry}[2]{%
  \noindent
  \begin{tabularx}{\linewidth}
    {@{}>{\raggedright\arraybackslash}X@{\hspace{6pt}}r@{}}
    \textbf{#1} & #2
  \end{tabularx}\par
}

\newcommand{\bullets}{
  \begin{itemize}[
    leftmargin=1.2em,
    topsep=2pt,
    itemsep=1pt,
    parsep=0pt,
    partopsep=0pt
  ]
}

\newcommand{\finishbullets}{
  \end{itemize}
}

\begin{document}

%----------HEADING----------
\begin{center}
    {\LARGE\bfseries Arumalla Mohit Krishna}\\[3pt]
    +1 (631) 352-8643 $|$
    \href{mailto:mohitkrishna.arumalla@stonybrook.edu}
    {mohitkrishna.arumalla@stonybrook.edu} $|$
    \href{https://linkedin.com/in/mohit-krishna-arumalla}
    {LinkedIn}
\end{center}
\vspace{-8pt}

%-----------EDUCATION-----------
\section{Education}

\entry
  {\textbf{Stony Brook University}}
  {\textbf{Aug 2025 -- Dec 2026}}

M.S. in Computer Science \hfill New York, USA\\
\textit{
Coursework: Natural Language Processing,
Introduction to Computer Vision, Vision Language Modeling
}
\vspace{3pt}

\entry
  {\textbf{Birla Institute of Technology and Science, Pilani}}
  {\textbf{Aug 2019 -- 2023}}

B.E. in Electronics \& Communication; GPA: 8.02/10 \hfill India
%-----------PROFESSIONAL SUMMARY-----------
\section{Professional Summary}

Software engineer with 2+ years of experience in \textbf{software development} and \textbf{data analysis}, currently pursuing an M.S. in Computer Science at Stony Brook University. Hands-on skills in \textbf{algorithm development}, \textbf{computational modeling}, and programming with strong interest in \textbf{computational physics} and research applications. Experienced in Python, Go, C++, and working with distributed data systems and infrastructure.

%-----------TECHNICAL SKILLS-----------
\section{Technical Skills}

\textbf{Languages:}
Python, Go, C++, SQL, Java\\

\textbf{Scientific \& Research Tools:}
NumPy, pandas, scikit-learn, PyTorch, TensorFlow/Keras, Monte Carlo methods, image processing, scientific computing\\

\textbf{Data \& Systems:}
gRPC, Kafka, Apache Flink, Debezium, PostgreSQL, Elasticsearch, Airflow, BigQuery\\

\textbf{Infrastructure:}
Docker, Kubernetes, Linux, Git, CI/CD, AWS

%-----------EXPERIENCE-----------
\section{Experience}

\entry
  {\textbf{Software Engineer II $|$ Infra.Market, Bangalore}}
  {\textbf{Apr 2024 -- Aug 2025}}

\bullets

  \item Built semantic product search using Amazon Titan text embeddings and Elasticsearch vector search, supporting natural-language queries and large-scale data processing.

  \item Engineered a CDC pipeline with Flink, Kafka, Debezium, PostgreSQL, and RocksDB, improving data synchronization latency to ~500 ms with exactly-once semantics.

  \item Deployed Istio/Envoy across microservices and integrated distributed tracing for real-time monitoring and debugging in production.

  \item Contributed to software infrastructure improvements and tooling automation to enhance system reliability and scalability.

\finishbullets

\entry
  {\textbf{Software Engineer I $|$ Infra.Market, Chennai}}
  {\textbf{Jul 2023 -- Mar 2024}}

\bullets

  \item Developed and maintained gRPC Search and Location microservices, handling 500K+ daily requests for product discovery and geolocation services.

  \item Fixed PostgreSQL connection leaks using Gorm ORM, reducing connection-pool exhaustion by 85\% and improving backend service stability.

  \item Built Airflow pipelines to migrate data from BigQuery to Elasticsearch with optimized queries and monitoring for data freshness.

  \item Supported data analysis and system integration efforts that enhanced existing software components and overall pipeline efficiency.

\finishbullets

%-----------PROJECTS-----------
\section{Projects}

\entry
  {GPU-Optimized Personalized Product Search}
  {Aug 2026 -- Oct 2026}

\textit{Python, PyTorch, BGE, NVIDIA Triton, Go, gRPC}

\bullets

  \item Developed two-tower retrieval model with \texttt{BAAI/bge-small-en-v1.5} embeddings and trainable user tower integrating query and interaction data.

  \item Built inference pipeline featuring precomputed product embeddings, approximate nearest-neighbor retrieval, and a Go/gRPC gateway; integrated NVIDIA Triton for GPU model serving.

  \item Benchmarking latency, throughput, GPU utilization, and cost for real-time retrieval under different traffic patterns; evaluating ranking metrics for search quality.

\finishbullets

\entry
  {Attention-Based Hallucination Detection in NMT}
  {Apr 2026 -- May 2026}

\textit{Python, PyTorch, mBART, Transformers, scikit-learn}

\bullets

  \item Built English-German hallucination-detection pipeline leveraging teacher-forced mBART attention features using WMT22/WMT23 quality estimation data.

  \item Extracted various attention metrics like entropy, max mass, variance, and grounding ratios across decoder layers and heads.

  \item Trained classifiers at token, sentence, and head levels achieving strong F1 and ROC-AUC scores; proposed hallucination-aware reranking methods.

\finishbullets

\entry
  {Prompt/LoRA for Self-Supervised Learning}
  {Ongoing}

\textit{Python, PyTorch, Vision Transformers}

\bullets

  \item Investigating domain-specific visual prompting and LoRA adaptation techniques in self-supervised ViT architectures to improve cross-domain transfer.

  \item Studying model performance and parameter efficiency with and without prompting modules under domain shifts.

\finishbullets

\end{document}